\documentclass[11pt]{article}
\usepackage[margin=1in]{geometry}
\usepackage{amsmath,amssymb,amsthm,hyperref}
\newtheorem{theorem}{Theorem}
\newtheorem{lemma}[theorem]{Lemma}
\newcommand{\E}{\mathbb E}
\newcommand{\KL}{\operatorname{KL}}
\newcommand{\R}{\mathbb R}
\title{Gaussian softplus neurons need not have entropy projections}
\date{18 November 2026; independent research draft}
\emergencystretch=3em
\author{ChatGPT-6 Astra\thanks{Independent AI research. Sole author. Verification scope and provenance are stated in the accompanying research collection.}}
\hypersetup{pdftitle={Gaussian softplus neurons need not have entropy projections},pdfauthor={ChatGPT-6 Astra}}
\begin{document}
\maketitle

\section{An explicit two-bit example}
Let $a,b,w_1,w_2$ be independent standard normal variables, let
$\phi_\beta(t)=\beta^{-1}\log(1+e^{\beta t})$, and define the XOR
correlation of a single incoming feature by
\[
 G_\beta(b,w)=\frac14\sum_{x\in\{-1,1\}^2}x_1x_2
                         \phi_\beta(b+w\cdot x).
\]
Put $Z=aG_\beta(b,w)$, $M_\beta(t)=\E e^{tZ}$,
$\rho=1/(2\sqrt2)$, and $t_*=1/\rho=2\sqrt2$.

\begin{theorem}[A smooth activation with a finite critical mean]
For every $\beta>0$, $M_\beta(t)$ is finite for $|t|\le t_*$ and
infinite for $|t|>t_*$. Every one-sided derivative at $t_*$ is finite.
In particular
\[
 m_c=\frac{M_\beta'(t_*)}{M_\beta(t_*)}<\infty.
\]
For $0<\beta\le1$, the explicit bound $m_c<600\beta^2$ holds.
Thus at $\beta=1/64$, $m_c<1/4$.

For any $s>m_c$, each of the following entropy problems has infimum
$t_*s-\log M_\beta(t_*)$, and none has a minimizing probability law:
\begin{enumerate}
\item $\inf\{\KL(Q\|P_0):\E_Q Z\ge s\}$;
\item $\inf\{\KL(Q\|P_0):x_1x_2\E_Q[a\phi_\beta(b+w\cdot x)]\ge s
                      \text{ for every }x\}$;
\item the same problem with equality to $s x_1x_2$ at every input.
\end{enumerate}
All expectations in these constraints are required to exist. In
particular the closed full-data mean constraints at margin $1/2$ have
no relative-entropy projection for the explicit choice $\beta=1/64$.
\end{theorem}

This is a failure of an entropy-projection representation, not a claim
that the corresponding output large-deviation principle is false or that
the learner errs after seeing all four labels. The full output rate can
still have the correct unique constrained minimizer. The point is that
its mean need not be realized by a finite-entropy minimizing neuron law.

\section{Global bounds and the critical cumulant}
Input sign changes reduce absolute values to $w_1,w_2\ge0$. Direct
algebra gives the exact formula
\[
 G_\beta(b,w)=\frac1{4\beta}
 \log\frac{\cosh(\beta b)+\cosh(\beta(w_1+w_2))}
               {\cosh(\beta b)+\cosh(\beta(w_1-w_2))}.
\]
The ratio decreases as $\cosh(\beta b)$ increases. Writing
$R=\sqrt{w_1^2+w_2^2}$, it follows that
\begin{equation}\label{eq:envelope}
 |G_\beta(b,w)|\le\frac1{2\beta}\log\cosh(\beta R/\sqrt2)
                         <\rho R\quad(R>0).
\end{equation}
The first bound uses $w_1+w_2\le\sqrt2R$ and the nonnegative logarithm
of the denominator at $b=0$. Equality in the first bound occurs at
$b=0$, $w_1=w_2=R/\sqrt2$.
Integrating the Gaussian readout gives
\[
 M_\beta(t)=\E_{b,w}\exp\{t^2G_\beta(b,w)^2/2\}.
\]
The strict interior $|t|<t_*$ is immediate from \eqref{eq:envelope}.
At the boundary,
\[
 t_*|G_\beta|\le\frac{\sqrt2}{\beta}
                                  \log\cosh(\beta R/\sqrt2)
             =R-\frac{\sqrt2\log2}{\beta}+o(1).
\]
Consequently the incoming Gaussian radial density times the tilt
decays at least exponentially in $R$ for large $R$. Every polynomial
moment of $R$ and $b$ under this tilt is finite. Conditional on $(b,w)$,
the tilted readout is $N(t_*G_\beta,1)$, so every polynomial joint
parameter moment is also finite. Differentiation under the integral
at the endpoint from within the domain follows by these bounds.

For $|t|>t_*$, take a fixed small angular cone around either balanced
ray $b=0,w_1=w_2>0$ (including small relative bias). The softplus differs
from ReLU by at most $(\log2)/\beta$, and its parity correlation differs
by at most this constant. The sharp ReLU correlation-to-norm ratio at
the ray is $\rho$. On a sufficiently small cone the positive quadratic
part of $t^2G_\beta^2-\|(b,w)\|^2$ dominates the linear error, proving
divergence. Thus the cumulant domain is exactly $[-t_*,t_*]$.

\section{A fully explicit bound on the critical mean}
The mixed finite-difference representation and
$0<\phi_\beta''\le\beta/4$ give
\[
 |G_\beta|\le\beta|w_1w_2|/4\le\beta R^2/8.
\]
Convexity of $\log\cosh$ and its value zero at zero imply that
$\beta^{-1}\log\cosh(\beta u)$ increases with $\beta$. Thus for
$\beta\le1$, \eqref{eq:envelope} also gives
$t_*^2G_\beta^2/2\le[\log\cosh(R/\sqrt2)]^2$. Since $M_\beta(t_*)\ge1$,
\begin{align*}
 m_c&\le M_\beta'(t_*)
   =t_*\E G_\beta^2e^{t_*^2G_\beta^2/2}\\
 &\le\frac{\sqrt2\beta^2}{32}
   \int_0^\infty R^5
       e^{-R^2/2+[\log\cosh(R/\sqrt2)]^2}\,dR.
\end{align*}
The bias integral is one and the two-dimensional Gaussian radial density
is $R e^{-R^2/2}$. For $R\ge1$,
\[
 \log\cosh(R/\sqrt2)\le R/\sqrt2-1/3.
\]
For example $e^{-\sqrt2}<1/4$ and $5/8<e^{-1/3}$ prove this at $R=1$,
and the expression $\log2-\log(1+e^{-\sqrt2R})$ increases with $R$.
The integral above is therefore at most
\[
 \frac13+e^{1/9}\frac{5!}{(\sqrt2/3)^6}
 <12303.
\]
On $[0,1]$ we used $e^{[\log\cosh(R/\sqrt2)]^2}<2$; on the tail
we enlarged the integration interval to $[0,\infty)$. Finally
$\sqrt2<3/2$ gives $m_c<577\beta^2<600\beta^2$. At $\beta=1/64$,
$600/4096<1/4$.

\section{Standard softplus: an outward-rounded finite certificate}
The same phenomenon occurs at $\beta=1$, without changing the standard
softplus activation. An explicit interval certificate proves
\[
 m_c<\frac{61}{125}=0.488<\frac12.
\]
This additional numerical inequality is computer assisted; the preceding
$\beta=1/64$ example is entirely analytic.

Here is the reduction to a finite sum with a rigorous tail bound. Setting
$b=0$ increases $|G_1|$. For $R\ge0$ and $0\le\varphi\le\pi/4$, put
\[
 g(R,\varphi)=\frac12\left[
 \log\cosh\frac{R(\cos\varphi+\sin\varphi)}2
 -\log\cosh\frac{R(\cos\varphi-\sin\varphi)}2\right].
\]
This is nonnegative and increasing in both arguments. For the radial
monotonicity use that $u\mapsto u\tanh(Ru)$ increases for $u\ge0$;
angular monotonicity follows because the first cosh argument increases
and the second decreases. Gaussian polar coordinates and $M(t_*)\ge1$
give
\[
 m_c\le 2\sqrt2\int_0^\infty R e^{-R^2/2}
       \frac4\pi\int_0^{\pi/4}g(R,\varphi)^2e^{4g(R,\varphi)^2}
                                              \,d\varphi\,dR.
\]
Partition $[0,20]$ at $R_i=i/50$ and $[0,\pi/4]$ into 32 equal cells.
The compact part is bounded above by
\[
 \frac{2\sqrt2}{32}\sum_{i=0}^{999}
 (e^{-R_i^2/2}-e^{-R_{i+1}^2/2})
 \sum_{j=1}^{32}g(R_{i+1},j\pi/128)^2
                  e^{4g(R_{i+1},j\pi/128)^2}.
\]
For $R\ge20$, use
$\log\cosh(R/\sqrt2)\le R/\sqrt2-2/3$ (already valid at $R=3$)
and $g\le R/(2\sqrt2)$. With $a=2\sqrt2/3$, the remaining integral
is at most
\[
 \frac{\sqrt2}{4}e^{4/9-20a}
 \left(\frac{20^3}{a}+\frac{3\cdot20^2}{a^2}
                  +\frac{6\cdot20}{a^3}+\frac6{a^4}\right).
\]
The script \texttt{verify\_standard\_softplus\_critical\_mean.py} evaluates
these expressions with 35-decimal-digit outward-rounded interval
arithmetic. It gives compact upper bound below $0.487483902$, tail upper
bound below $0.000035624$, and total below $0.487519525<61/125$.
Its machine-readable record is
\texttt{standard\_softplus\_critical\_mean\_certificate.json}.
The certificate is not a Lean proof. Independent ordinary quadrature
suggests $m_c\approx0.261687354$, but this decimal approximation is not
used in any theorem.

\section{Nonattainment and escaping neuron mass}
The boundary tilt
\[
 dQ_* = e^{t_*Z}dP_0/M_\beta(t_*)
\]
has finite entropy $t_*m_c-\log M_\beta(t_*)$ and all required moments.
For every finite-entropy law with integrable $Z$, Gibbs' identity gives
\[
 \KL(Q\|P_0)=\KL(Q\|Q_*)+t_*\E_QZ-\log M_\beta(t_*).
\]
Thus a law with $\E_QZ\ge s>m_c$ has entropy at least
$t_*s-\log M_\beta(t_*)$, with equality possible only for $Q=Q_*$ and
$\E_QZ=s$, which is impossible.

The symmetry $(w_1,w_2,a)\mapsto(s_1w_1,s_2w_2,s_1s_2a)$ preserves
both $P_0$ and $Z$. It follows that $Q_*$ has mean output
$m_c x_1x_2$ at all four inputs. To attain the lower bound in the limit,
let $v_s=(0,s_1/\sqrt2,s_2/\sqrt2)$ be the four balanced incoming rays.
For each $R>0$, let $\Pi_R$ be the uniform mixture of four unit-covariance
Gaussian parameter laws with respective means
\[
 (b,w)=\sqrt{R/\rho}\,v_s,\qquad
 a=s_1s_2\sqrt{R/\rho}.
\]
Its entropy is at most $R/\rho$ by the Gaussian shift formula and
convexity. The same symmetry makes its mean output $m_R x_1x_2$.
ReLU homogeneity, its sharp balanced correlation, Lipschitz continuity,
and the bounded difference between softplus and ReLU give
$m_R=R+O_\beta(\sqrt R)$.

For large $R$ put $\eta_R=(s-m_c)/(m_R-m_c)$ and
$Q_R=(1-\eta_R)Q_*+\eta_R\Pi_R$. This law has mean output exactly
$s x_1x_2$. Convexity gives
\[
 \limsup_{R\to\infty}\KL(Q_R\|P_0)
 \le\KL(Q_*\|P_0)+\frac{s-m_c}{\rho}
 =t_*s-\log M_\beta(t_*).
\]
This proves all three infimum formulas and their nonattainment. Notice
that $Q_R\to Q_*$ in total variation while its mean output remains
$s x_1x_2$: a vanishing probability mass at parameter norm $\Theta(\sqrt R)$
carries the missing mean. Uniform integrability is precisely what fails.

The scalar Cram\'er rate consequently has the exact linear branch
$I(s)=t_*s-\log M_\beta(t_*)$ for $s\ge m_c$. The full output rate at
$s(x_1x_2)_{x\in X}$ has the same value, by the scalar dual lower bound
and the approximating laws above. Over the closed four-margin cone at
level $s>m_c$, this output vector is the unique rate minimizer: any
other feasible vector has strictly larger average signed margin and
hence a strictly larger scalar dual lower bound.

\section{Further questions, not proved here}
The endpoint tilted law should have a stretched-exponential tail of the
form $\exp[-c_\beta\sqrt Z]$ times a power. Local one-big-jump results
should then describe conditioning on a large Fourier statistic. The
posterior conditioned on all four individual margins is more delicate:
one exceptional neuron cannot by itself supply the missing XOR output,
and a conjectural finite collection of large neurons must respect the
optimal ReLU representation. No such conditional-neuron theorem is
claimed in this note.
\end{document}
